\section*{Chapter \ref{ch:mx:circuit-breakers-and-halts} --- Circuit Breakers and Halts}

\begin{solution}{exo:mx:circuit-breakers-and-halts:1}
$99.93\times0.9985=99.780$ rounded down to 99.78, and $99.93\times1.0015=100.080$ rounded up to 100.08.
\end{solution}

\begin{solution}{exo:mx:circuit-breakers-and-halts:2}
15 minutes: the threshold was crossed at 9:34:13 a.m. and the \omterm{def:mx:circuit-breakers-and-halts:pause}{reopening auctions} began at 9:49:13 a.m. A Level~2 decline at or after 3:25 p.m. would not
have halted the market.
\end{solution}

\begin{solution}{exo:mx:circuit-breakers-and-halts:3}
43.9\% with no mechanism, 41.1\% with the band, 29.4\% with \omterm{def:mx:circuit-breakers-and-halts:velocity}{velocity logic}, 23.2\% with the market-wide halt.
\end{solution}

\begin{solution}{exo:mx:circuit-breakers-and-halts:4}
Its lower edge is a barrier: bids below it are rejected and market sell orders stop at it, their remainder cancelled. The selling that would have walked
the book further never trades, so the price cannot fall past the band while the band itself only follows the minute's average price down.
\end{solution}

\begin{solution}{exo:mx:circuit-breakers-and-halts:5}
Its pauses are short and each \omterm{def:mx:circuit-breakers-and-halts:pause}{reopening auction} meets the same one-sided selling: the pause stops trading without bringing new buyers, and the fall
resumes; meanwhile the seller's orders sent during the pause wait for the reopening and add to its imbalance.
\end{solution}

\begin{solution}{exo:mx:circuit-breakers-and-halts:6}
Condition on the distance to the limit, measure the speed of approach or the probability of reaching the limit, and compare with a control that
lacks the mechanism (another market, a period before the rule, a limit that was not in force); separate the mechanical barrier (orders refused at the
limit) from behaviour (orders advanced in time), for example with order-level data on who sends orders near the limit.
\end{solution}

\begin{solution}{exo:mx:circuit-breakers-and-halts:7}
At 3 ticks: magnet $0.727-0.694=+0.03$ with a standard error of about 0.09; barrier $0.694-0.816=-0.12$ with about 0.09. Neither is beyond two
standard errors at one distance; the barrier's sign is the same at ten of the eleven distances.
\end{solution}

\begin{solution}{exo:mx:circuit-breakers-and-halts:8}
Fewer stocks reach their limits may mean the band blocks trades at its edge (the barrier effect), not that the market is calmer; volume and fills refused
at the band, and the prices after it, say whether risk was reduced or postponed. A before-after comparison also ignores what else changed.
\end{solution}

\begin{solution}{pb:mx:circuit-breakers-and-halts:1}
\textbf{1.} See the definitions.
\textbf{2.} Bands against runaway or erroneous prices, pauses against moves faster than liquidity returns, halts against undigested news, limits all day.
\textbf{3.} Through the engine's controls: R sets a band, P changes the phase; the policy checks every half second.
\textbf{4.} 7\%, 13\% (15 minutes, before 3:25 p.m.) and 20\% (rest of the day); Level~1 halts on 9, 12, 16 and 18 March 2020.
\textbf{5.} 40\,000 shares sold as sixty \omterm{def:mx:the-limit-order-book:orders}{market orders} over a minute; no mechanism, a dynamic band with pauses, \omterm{def:mx:circuit-breakers-and-halts:velocity}{velocity logic}, a market-wide halt.
\textbf{6.} Falls of 27.4, 20.0, 24.9 and 18.9 ticks; 0, 0, 79 and 48 seconds halted; 17\,580, 16\,440, 11\,752 and 9\,279 shares filled; costs 12.2, 9.6,
13.2 and 6.6 ticks.
\textbf{7.} It stops all trading at the first level and takes the most orders out of the market.
\textbf{8.} A pause after a move of a set number of ticks within a set time; orders may be entered, modified and cancelled in the reserved state.
\textbf{9.} Circuit breakers may make agents advance trades and so increase variability and exacerbate moves.
\textbf{10.} Sellers who, when a band is enforced, sell faster the closer the best bid is to its lower edge.
\textbf{11.} For each session and distance, the first time within $x$ ticks and whether the band is reached within 30 seconds; band with anticipators,
band alone, no band.
\textbf{12.} \emph{Named result}: at 5 ticks from the lower band, the probability of reaching it within 30 seconds is 0.66 with the band and
anticipating sellers, 0.63 with the band alone and 0.75 without a band; the simulated magnet effect averages $+0.03$ over 3 to 7 ticks (not
distinguishable from zero in 80 sessions), the band's barrier effect $-0.12$.
\textbf{13.} It removes the sellers' orders at its edge, which outweighs the extra selling it invites.
\textbf{14.} Indicative prices during the call, collars on the reopening price, extensions when imbalance is large, and enough time for the other side
to arrive.
\textbf{15.} London single-stock circuit breakers improved the liquidity and reduced the volatility of the stocks that kept trading.
\textbf{16.} They ended before new buyers arrived.
\textbf{17.} More anticipating traders, or a band that pauses instead of refusing trades.
\textbf{18.} Protecting prices by refusing trades: shorter falls, fewer shares sold.
\textbf{19.} That the bands and halts limit falls mainly by refusing trades, that any magnet effect is small next to that, and that a pause helps only if
liquidity can arrive before it ends.
\textbf{20.} Time for liquidity to return, paid for with trades that do not happen.
\end{solution}

\begin{solution}{iq:mx:circuit-breakers-and-halts:1}
Stop the algorithm's clock, decide whether to take part in the \omterm{def:mx:circuit-breakers-and-halts:pause}{reopening auction} (and at what limit, given the indicative price), and re-plan the rest
of the order for after the reopening; do not send \omterm{def:mx:the-limit-order-book:orders}{market orders} into the call.
\iqlookfor{Participation in the reopening; re-planning.}
\end{solution}

\begin{solution}{iq:mx:circuit-breakers-and-halts:2}
The acceleration of prices toward a limit because traders rush before it binds; test by conditioning on the distance to the limit and comparing with a
market or period without it, separating the barrier from behaviour.
\iqlookfor{Conditioning on distance; a control; mechanics against behaviour.}
\end{solution}

\begin{solution}{iq:mx:circuit-breakers-and-halts:3}
Cash equities and options stop but futures may keep trading within their own limits, so hedges can be adjusted only where markets are open; risk
limits must use the last prices with care, and margin calls may come before the reopening.
\iqlookfor{Cross-market differences; stale prices.}
\end{solution}

\begin{solution}{iq:mx:circuit-breakers-and-halts:4}
A halt: stop and wait for the reopening, cancel or hold orders per policy; a pause: participate in the \omterm{def:mx:circuit-breakers-and-halts:pause}{reopening auction} or wait; a band rejection:
re-price within the band or wait, never retry blindly; log and alert in every case.
\iqlookfor{State-aware handling; no blind retries.}
\end{solution}

\begin{solution}{iq:mx:circuit-breakers-and-halts:5}
\omterm{def:mx:circuit-breakers-and-halts:velocity}{Velocity logic} reacts to speed (ticks per second) with a pause of seconds; limit up-limit down reacts to distance from a moving reference with a
limit state and a five-minute pause: one is about fast moves, the other about large ones.
\iqlookfor{Speed against distance; pause lengths.}
\end{solution}

\begin{solution}{iq:mx:circuit-breakers-and-halts:6}
Evidence on how often bands bind and pauses trigger, on price continuation after reopenings, and on spillovers to other stocks; expect fewer triggers
and less refused volume, but larger moves inside the bands and less time for liquidity to return.
\iqlookfor{Trade-offs; the measures to watch.}
\end{solution}
